\documentclass[../main]{subfiles}
\begin{document}
\ifSubfilesClassLoaded{\setcounter{page}{4}}{}
\section{Conclusione}
\label{sec:11_conclusione}

Solo ora, dopo aver delineato i principi fondamentali del comunismo e averne individuato i semi, da cui esso può svilupparsi come sistema sociale, ci sentiamo legittimati a passare alle definizioni. Altrimenti, queste risulterebbero mere astrazioni, prive di significato, e non costituirebbero una risposta alla domanda che abbiamo posto all'inizio: «Che cos'è il comunismo?».

Il comunismo (la società comunista) è un'associazione libera di individui che organizzano razionalmente la propria esistenza, ponendo al centro, in modo consapevole, la propria libera attività.

Il comunismo è una forma di libera attività intrinsecamente collettiva, la forma di una società che si crea e determina il proprio sviluppo, divenendo essa stessa totalità.

Come affermarono Marx ed Engels nel «Manifesto», il comunismo è «un'associazione in cui il libero sviluppo di ciascuno è condizione per il libero sviluppo di tutti».

È tempo di riaffermare, con rinnovato vigore, un motto che non ha perduto la sua forza:

PROLETARI DI TUTTI I PAESI, UNITEVI!
\end{document}
